\chapter{Bitfields\label{chap:Bitfields}}

\hypertarget{bitfieldterm}{}
Bitvector types allow defining named slices, which we refer to as \emph{bitfields}.
Bitfields can be read from and written to.

\ExampleDef{A bitvector type with bitfields}
\listingref{bitfields1} declares a global variable whose type is a bitvector with bitfields.
\ASLListing{A bitvector type with bitfields}{bitfields1}{\definitiontests/Bitfields.asl}
\begin{itemize}
  \item The expression \texttt{myData.flag} evaluates to the same value as \verb|myData[4]|, with type \verb|bits(1)|;
  \item The expression \texttt{myData.data} evaluates to the same value as\\ \verb|myData[3:0] :: myData[8:5]|, with type \verb|bits(8)|;
  \item There is no bitfield which accesses \verb|myData[15:10]|;
  \item The \verb|value| field overlaps with the other fields;
  \item The slices \verb|3:0| and \verb|8:5| which define \texttt{data} do not overlap each other,
        but they do overlap \verb|value|.
\end{itemize}
Note that in the \texttt{data} bitfield, bits |3:0| come before and are more significant than
bits |8:5|, which is a different order from their occurrence in \texttt{myData}.

\hypertarget{def-singleslice}{}
\hypertarget{def-rangeslice}{}
\hypertarget{def-lengthslice}{}
\hypertarget{def-scaledslice}{}
We refer to a slice of the form \texttt{[$\ve$]} as a \singleslice,
a slice of the form \texttt{[$\veone$:$\vetwo$]} as a \rangeslice,
a slice of the form \texttt{[$\veone$+:$\vetwo$]} as a \lengthslice,
and a slice of the form \texttt{[$\veone$*:$\vetwo$]} as a \scaledslice.

\ChapterOutline
\begin{itemize}
  \item \secref{Nested Bitfields} describes nested bitfields;
  \item \FormalRelationsRef{Bitfields} defines the formal relations for bitfields;
  \item \SyntaxRef{Bitfields} defines the syntax of bitfields;
  \item \AbstractSyntaxRef{Bitfields} defines the abstract syntax of bitfields;
  \item \TypeRulesRef{Bitfields} defines the type system rules for bitfields; and
  \item \SemanticsRulesRef{Bitfields} defines the dynamic semantics of bitfields.
\end{itemize}

\section{Nested Bitfields\label{sec:Nested Bitfields}}
Bitfields may have nested bitfields. This can have several uses, one of which is being able to define two
different views of a register.

This section defines several properties of bitfields, illustrated by the example below.
These are used to define \RequirementRef{BitfieldAlignment}:
\begin{itemize}
\item The \hypertarget{def-absolutename}{\absolutename} of a (possibly-nested) bitfield is the list of identifiers
that start from the
name of the top-level bitfield containing it and follows with the names of the nested bitfields,
until the name of the bitfield at hand. We denote an absolute name by a period-separated list of names,
similar to the expression referring to it;
\item The \hypertarget{def-absoluteslice}{\absoluteslice} of a bitfield is comprised of the slices it defines, with respect to its containing
bitvector type;
\item The \hypertarget{def-absolutebitfield}{\absolutebitfield} of a given bitfield consists of absolute name and its
\absoluteslices; and
\item The \hypertarget{def-bitfieldscope}{\bitfieldscope} of a bitfield is its \absolutename, with the last
identifier (the bitfield's name) removed.
For example, the \bitfieldscope{} of a bitfield with the absolute name \texttt{fmt0.later1.remainder.moving}
is \texttt{fmt0.later1.remainder} and the \bitfieldscope{} of a bitfield with the absolute name
\texttt{fmt0} is the empty list.
\end{itemize}

\begin{definition}[Common Bitfield Scope]
We say that two bitfields exist in the same \bitfieldscope{} if the \bitfieldscope{} of one is a prefix
of the \bitfieldscope{} of the other.
\end{definition}

\ExampleDef{Bitfields in the Same Scope}
The following table shows pairs of bitfields given by their \absolutenames{}.
The pair in the left column exist in the same \bitfieldscope{},
since \texttt{fmt0.later1} is a prefix of
\texttt{fmt0.later1.remainder}.
In contrast, the pair in the right column do not,
as neither \texttt{fmt0} nor \texttt{fmt1} is a prefix of the other.

\begin{center}
\begin{tabular}{ll}
\hline
\textbf{Same scope} & \textbf{Different scopes}\\
\hline
\texttt{fmt0.later1.remainder} & \texttt{fmt0.common} \\
\texttt{fmt0.later1.remainder.moving} & \texttt{fmt1.common}\\
\hline
\end{tabular}
\end{center}

\RequirementDef{BitfieldAlignment}
For every two bitfields of a given bitvector type,
if they share the same name and exist in the same \bitfieldscope\ then their \absoluteslices\ must be
equal.
This is formalized in \TypingRuleRef{CheckCommonBitfieldsAlign} and exemplified in \listingref{nestedbitfields}.

\ExampleDef{A bitvector type with nested bitfields}
\listingref{nestedbitfields} shows a bitvector type with nested bitfields, along with the absolute bitfield
defined for each bitfield.
\ASLListing{A bitvector type with nested bitfields}{nestedbitfields}{\definitiontests/Bitfields_nested.asl}

%%%%%%%%%%%%%%%%%%%%%%%%%%%%%%%%%%%%%%%%%%%%%%%%%%%%%%%%%%%%%%%%%%%%%%%%%%%%%%%%
\FormalRelationsDef{Bitfields}
\paragraph{Syntax:} Bitfields are grammatically derived from $\Nbitfield$
  and lists of bitfields are grammatically derived from $\Nbitfields$.

\paragraph{Abstract Syntax:} Bitfields are derived in the abstract syntax from
  $\bitfield$ and generated by $\buildbitfield$ (see \ASTRuleRef{Bitfield}).

\paragraph{Typing:} Bitfields are annotated by $\annotatebitfields$.

\paragraph{Semantics:} Bitfields are not directly associated with a dynamic semantics.
  Their semantics are given as part of slicing expressions
  by $\evalexpr$ (see \secref{SlicingExpressions}) and
  $\evallexpr$ (see \secref{BitvectorSliceAssignmentExpressions}).

%%%%%%%%%%%%%%%%%%%%%%%%%%%%%%%%%%%%%%%%%%%%%%%%%%%%%%%%%%%%%%%%%%%%%%%%%%%%%%%%
\SyntaxDef{Bitfields}
\begin{flalign*}
\Nbitfields \derives \ & \Tlbrace \parsesep \TClistZero{\Nbitfield} \parsesep \Trbrace &\\
\Nbitfield \derives \ & \Nslices \parsesep \Tidentifier &\\
                  |\ & \Nslices \parsesep \Tidentifier \parsesep \Nbitfields &\\
                  |\ & \Nslices \parsesep \Tidentifier \parsesep \Tcolon \parsesep \Nty &\\
\end{flalign*}

%%%%%%%%%%%%%%%%%%%%%%%%%%%%%%%%%%%%%%%%%%%%%%%%%%%%%%%%%%%%%%%%%%%%%%%%%%%%%%%%
\AbstractSyntaxDef{Bitfields}
\RenderType[remove_hypertargets]{bitfield}

\ASTRuleDef{Bitfields}
\RenderRelation{build_bitfields}


\begin{mathpar}
\inferrule{
  \buildtclist[\buildbitfield](\vbitfields) \astarrow \vbitfieldasts
}{
  {
    \begin{array}{r}
  \buildbitfields(\Nbitfields(\Tlbrace, \namednode{\vbitfields}{\TClistZero{\Nbitfield}}, \Trbrace)) \astarrow\\
  \overname{\vbitfieldasts}{\vastnode}
  \end{array}
  }
}
\end{mathpar}

\ASTRuleDef{Bitfield}
\RenderRelation{build_bitfield}


\begin{mathpar}
\inferrule[simple]{}{
  \buildbitfield(\Nbitfield(\punnode{\Nslices}, \Tidentifier(\vx))) \astarrow
  \overname{\BitFieldSimple(\vx, \astof{\vslices})}{\vastnode}
}
\end{mathpar}

\begin{mathpar}
\inferrule[nested]{}{
  {
    \begin{array}{r}
      \buildbitfield(\Nbitfield(\punnode{\Nslices}, \Tidentifier(\vx), \punnode{\Nbitfields})) \astarrow\\
      \overname{\BitFieldNested(\vx, \astof{\vslices}, \astof{\vbitfields})}{\vastnode}
    \end{array}
  }
}
\end{mathpar}

\begin{mathpar}
\inferrule[type]{}{
  {
    \begin{array}{r}
      \buildbitfield(\Nbitfield(\punnode{\Nslices}, \Tidentifier(\vx), \Tcolon, \punnode{\Nty})) \astarrow\\
      \overname{\BitFieldType(\vx, \astof{\vslices}, \astof{\tty})}{\vastnode}
    \end{array}
  }
}
\end{mathpar}

%%%%%%%%%%%%%%%%%%%%%%%%%%%%%%%%%%%%%%%%%%%%%%%%%%%%%%%%%%%%%%%%%%%%%%%%%%%%%%%%
\TypeRulesDef{Bitfields}
\TypingRuleDef{TBitFields}
\RenderRelation{annotate_bitfields}

\ExampleDef{Typing Bitfields}
The bitfields declared in \listingref{bitfields1} are well-typed and their total width
is $16$.

\RenderProseAndFormally{annotate_bitfields}
\CodeSubsection{\TBitFieldsBegin}{\TBitFieldsEnd}{../Typing.ml}

\TypingRuleDef{BitFieldGetName}
\RenderRelation{bitfield_get_name}

\ExampleDef{Bitfield Names}
In \listingref{nestedbitfields}, the names of bitfields are:
\verb|fmt|, \verb|common|, \verb|layer1|, \verb|remainder|, \verb|moving|,
\verb|extra|, and \verb|fmt1|.

\RenderProseAndFormally{bitfield_get_name}

\TypingRuleDef{BitFieldGetSlices}
\RenderRelation{bitfield_get_slices}

\ExampleDef{The Slices of a Bitfield}
In \listingref{nestedbitfields},
the slices associated with the bitfield \verb|layer1| are \verb|4:13, 12:2, 1, 0|.


\RenderProseAndFormally{bitfield_get_slices}

\TypingRuleDef{BitFieldGetNested}
\RenderRelation{bitfield_get_nested}

\ExampleDef{The Bitfields Nested in a Bitfield}
In \listingref{nestedbitfields},
the bitfields nested in the bitfield \verb|layer1| consist of the singleton list \verb|remainder|,
which does not include \verb|moving| and \verb|extra| -- those are nested in \verb|remainder|.
The bitfields nested in the bitfield \verb|fmt| make up an empty list.

\RenderProseAndFormally{bitfield_get_nested}

\TypingRuleDef{TBitField}
\RenderRelation{annotate_bitfield}

\ExampleDef{Well-typed Bitfields}
In \listingref{welltypedbitvectortypes}, all the uses of bitvector types with bitfields are well-typed.
\ASLListing{Well-typed bitfields}{welltypedbitvectortypes}{\typingtests/TypingRule.TBitField.asl}


\RenderProseAndFormally{annotate_bitfield}
\CodeSubsection{\TBitFieldBegin}{\TBitFieldEnd}{../Typing.ml}

\TypingRuleDef{CheckSlicesInWidth}
\RenderRelation{check_slices_in_width}

See
\ExampleRef{Converting Disjoint Slices to Positions} and
\ExampleRef{Checking Whether Slice Positions Fit in a Bitvector Width}.


\RenderProseAndFormally{check_slices_in_width}
\CodeSubsection{\CheckSlicesInWidthBegin}{\CheckSlicesInWidthEnd}{../Typing.ml}

\TypingRuleDef{CheckPositionsInWidth}
\RenderRelation{check_positions_in_width}

\ExampleDef{Checking Whether Slice Positions Fit in a Bitvector Width}
In \listingref{bitfieldslicetopositions}, all slices declared for all bitfields
fit in the bitvector width $16$, whereas the positions defined for the bitfield
\verb|value| in \listingref{checkpositionsinwidth} exceed the bitvector width $16$.

\ASLListing{An Out of Width Slice}{checkpositionsinwidth}{\typingtests/TypingRule.CheckPositionsInWidth.bad.asl}


\RenderProseAndFormally{check_positions_in_width}

\TypingRuleDef{DisjointSlicesToPositions}
\RenderRelation{disjoint_slices_to_positions}

This function assumes that all expressions in $\vslices$ are \pureterm{} and thus deterministic.

\ExampleDef{Converting Disjoint Slices to Positions}
In \listingref{bitfieldslicetopositions}, the slices \verb|3:0| and \verb|5+:3|, declared for the bitfield \verb|data|
yield the set of positions $\{0, 1, 2, 3, 5, 6, 7\}$.
Whereas the slices \verb|3:0| and \verb|5:3|, declared for the bitfield \verb|data|
in \listingref{disjointslicestopositions} overlap, since they have $3$ in common.

\ASLListing{Overlapping Slices}{disjointslicestopositions}{\typingtests/TypingRule.DisjointSlicesToPositions.bad.asl}


\RenderProseAndFormally{disjoint_slices_to_positions}

\TypingRuleDef{BitfieldSliceToPositions}
\RenderRelation{bitfield_slice_to_positions}

The function assumes that AST nodes labelled with $\SliceSingle$, $\SliceRange$, and $\SliceStar$ have been reduced
to $\SliceLength$ (via $\annotateslices$), and that all expressions in $\vslice$ are \pureterm{} and thus deterministic.

\ExampleDef{Converting Bitfield Slices to Positions}
\taref{bitfieldslicetopositions} shows the optional set of positions associated with each slice
in \listingref{bitfieldslicetopositions}, followed by examples of erroneous slices.

\begin{table}[h]
\caption{Converting Bitfield Slices to Positions\label{ta:bitfieldslicetopositions}}
\begin{center}
\begin{tabular}{ll}
  \textbf{Slice} & \textbf{Optional Set of Positions}\\
  \hline
  \verb|4|    & $\some{\{4\}}$ \\
  \verb|3:0|  & $\some{\{0, 1, 2, 3\}}$\\
  \verb|5+:3| & $\some{\{5, 6, 7\}}$\\
  \verb|3*:4| & $\some{\{12, 13, 14, 15\}}$\\
  \hline
  \verb|0:3|  & $\BadSlices$\\
  \verb|5+:0| & $\BadSlices$\\
  \verb|4*:0| & $\BadSlices$\\
\end{tabular}
\end{center}
\end{table}

\ASLListing{Converting bitfield slices to positions}{bitfieldslicetopositions}{\typingtests/TypingRule.BitfieldSliceToPositions.asl}


\RenderProseAndFormally{bitfield_slice_to_positions}

\TypingRuleDef{EvalSliceExpr}
\RenderRelation{eval_slice_expr}

\ExampleDef{Evaluating Expressions in Slices}
The specification in \listingref{EvalSliceExpr} shows two kinds of slices:
the slice \\
\verb|static_func{4}(4):0| is used to define a bitfield (\verb|data|)
and therefore must be \staticallyevaluableterm,
whereas the slice \verb|N DIV 2:0| is used in an \assignableexpression, which means
it need not be \staticallyevaluableterm{} (it is indeed not, in this example).
\ASLListing{Expressions in slices}{EvalSliceExpr}{\typingtests/TypingRule.EvalSliceExpr.asl}


\RenderProseAndFormally{eval_slice_expr}

\TypingRuleDef{CheckCommonBitfieldsAlign}
\RenderRelation{check_common_bitfields_align}

\TypingRuleRef{TBits} (transitively) guarantees that $\vwidth > 0$ holds.

We represent \absolutebitfields\ by the type
\RenderType{TAbsField}
where the component $\name$ is a list of identifiers corresponding to the \absolutename,
and the component $\vslice$ corresponds to an \absoluteslice{} by listing the indices
into the containing bitvector type.\footnote{An implementation of the type system may compactly represent the list of indices
via a list of intervals, each represented by its limits.}

\ExampleDef{Ill-typed Bitfields}
\listingref{CommonBitfieldsAlignError} shows an example where the two bitfields named \texttt{common}
exist in the same \bitfieldscope, but their \absoluteslices\ are not the same.
Specifically, the \absoluteslice\ for the bitfield \texttt{common} is \texttt{[1:0]}
whereas the \absoluteslice\ for the bitfield \texttt{sub.common} is \texttt{[0, 1]}.
Typechecking this example results in the \typingerrorterm{} \BadSlices.

\ASLListing{An example where two bitfields of the same name (\texttt{common}) exist in the same scope but have different absolute slices}{CommonBitfieldsAlignError}{\typingtests/TypingRule.CheckCommonBitfieldsAlign.Error.asl}

\RenderProseAndFormally{check_common_bitfields_align}

\TypingRuleDef{BitfieldsToAbsolute}
\RenderRelation{bitfields_to_absolute}

See \ExampleRef{A bitvector type with nested bitfields}.


\RenderProseAndFormally{bitfields_to_absolute}

\TypingRuleDef{BitfieldToAbsolute}
\RenderRelation{bitfield_to_absolute}

See \ExampleRef{A bitvector type with nested bitfields}.


\RenderProseAndFormally{bitfield_to_absolute}

\TypingRuleDef{SelectIndicesBySlices}
\RenderRelation{select_indices_by_slices}

% The reference abstracts away from the implementation,
% which represents sequence of indices via intervals.

\ExampleDef{Selecting from a List of Indices}
The following are some examples of selecting indices:
\[
\begin{array}{lcl}
\selectindicesbyslices([13,1,6,4,9], [0,1,2,3,4]) &\typearrow& [13,1,6,4,9]\\
\selectindicesbyslices([13,1,6,4,9], [4,3,2,1,0]) &\typearrow& [9,4,6,1,13]\\
\selectindicesbyslices([13,1,6,4,9], [0,4,2,3])   &\typearrow& [13,9,6,4]\\
\end{array}
\]

\RenderProseAndFormally{select_indices_by_slices}

\TypingRuleDef{AbsoluteBitfieldsAlign}
\RenderRelation{absolute_bitfields_align}

See \ExampleRef{A bitvector type with nested bitfields} where all \absolutebitfields{} align
and \ExampleRef{Ill-typed Bitfields} where not all \absolutebitfields{} align.


\RenderProseAndFormally{absolute_bitfields_align}

\TypingRuleDef{SliceToIndices}
\RenderRelation{slice_to_indices}

\ExampleDef{Converting a Slice to a List of Indices}
The list of indices for the slice \\
$\SliceLength(\AbbrevEBinop{\ADD}{\ELInt{1}}{\ELInt{4}}, \ELInt{3})$ is
$[7, 6, 5]$.


\RenderProseAndFormally{slice_to_indices}

\SemanticsRulesDef{Bitfields}
The type system annotates expressions that read from a bitfield into corresponding slicing \rhsexpressions{}.
Therefore, reading from a bitfield is done by evaluating a corresponding slice expression (see \SemanticsRuleRef{ESlice}).

The type system annotates expressions that write to a bitfield into corresponding slicing \assignableexpressions{}.
Therefore, writing to a bitfield is done by evaluating a corresponding assignable slice expression (see \SemanticsRuleRef{LESlice}).
